\section{Introduction}
\label{sec:intro}

An LLM agent's behavior is set as much by its harness as by its model. The harness is the software
layer that assembles the model's input at every step, executes the actions the model selects, and
decides whether to continue; it also holds state, verifies output, bounds cost, isolates execution,
and records what happened. Hold the model fixed and change the harness, and the reported score moves
\citep{guo2026qa2task}, so an agent comparison that does not disclose its harness cannot be
interpreted \citep{zhang2026harnessreporting}. The design question is how harnesses are built and
which design choices change outcomes.

Seven prior surveys and taxonomies address harnesses, and none can answer that question.
\citet{guo2026qa2task} decompose the harness into six runtime responsibilities and alone compile
same-model, different-harness leaderboard evidence, but the harness enters that evidence as a name
rather than as coded design. \citet{rombaut2026insidescaffold} reads the source of thirteen coding
agents with a file and line behind every claim, and \citet{barbaste2026anatomy} read the source of
eleven harnesses at pinned releases.
\citet{li2026harnessengineering}, \citet{meng2026harnesssurvey}, \citet{zhong2026responsibilities},
and \citet{harnesscategorical2026} contribute a layered taxonomy, a formal component tuple, a maturity
ladder, and a categorical semantics. Four technical limits are shared. First, none samples from an
enumerated frame of released systems, so none can estimate how far its well-known systems depart from
the population. Second, none codes design with re-openable evidence per cell beyond thirteen systems,
so a claimed value cannot be rechecked against its source. Third, none was registered before coding
or reports agreement per dimension, and \citet{li2026harnessengineering} state that they do not report
Cohen's kappa; without a reliability gate, consistent use of a category is assumed rather than shown.
Fourth, none records documented silence as a state of its own, so where the evidence is documents
rather than code, an undocumented component and an absent one receive the same mark. A catalog built
this way says what exists and what to call it; it cannot measure what the field discloses or relate
design to outcome.

We address the design question with a review registered on OSF (\texttt{osf.io/ab2wn}) before coding
and reported under PRISMA 2020 \citep{page2021prisma}, built from three frameworks
(Figure~\ref{fig:pipeline}). \emph{Two-Tier Escalation Screening} settles each retrieved record at
title and abstract where it can and escalates the rest to a full-text reading of paper, repository,
and documentation together. Grouping the included records into systems gives a sampling frame of
6,504 harnesses, sampled by stratum and weighted back. \emph{Evidence-Anchored Coding} codes 1,256
systems on 38 dimensions in nine layers with a verbatim quote and a re-openable locator behind every
claimed value. Before the schema froze, each dimension was tested against a floor of
$\kappa \ge 0.6$ on a double-coded sample. The \emph{Three-State Cell Contract} gives every cell
exactly one state: a value, \texttt{not\_reported} when the sources were read and are silent, or
\texttt{unresolved} when the cell could not be settled, so silence is counted and never rounded to
absence. This is not the first survey of agent harnesses. To our knowledge, it is the first
pre-registered review of the area to report per-dimension coding reliability, and the first to analyze
reported outcomes against coded design dimensions rather than harness identity.

\begin{figure}[t]
  \centering
  \includegraphics[width=\linewidth]{pipeline_overview}
  \caption{The study at a glance. Two-Tier Escalation Screening takes 27,747 screened records to a
  frame of 6,504 systems; Evidence-Anchored Coding turns 1,256 of them into 47,728 cells under the
  Three-State Cell Contract.}
  \label{fig:pipeline}
\end{figure}

The results say that the literature describes what a harness is and not what bounds it, and that its
outcome evidence cannot yet say which design choices matter. Of the 47,728 cells, 48.9\% record
documented silence, and the silence concentrates in the layers that govern safe deployment. Beneath a
comparability ceiling of 53 of 1,256 coded systems (4.2\%), of the two design estimates that survive, one is an
upper bound and the other is biased in both directions.

Four registered questions organize the paper.

\begin{enumerate}
  \item[\textbf{RQ1}] What are a harness's structural components, and can a taxonomy of them be applied
  at $\kappa \ge 0.6$ per dimension?
  \item[\textbf{RQ2}] How do design choices vary across systems, and how did that variation move
  between 2022 and 2026?
  \item[\textbf{RQ3}] Which design choices are associated with reported success, cost, or robustness,
  with benchmark and base model held constant?
  \item[\textbf{RQ4}] Which dimensions does the literature systematically fail to report?
\end{enumerate}

Our contributions follow.

\begin{enumerate}
  \item \textbf{An evidence-anchored release (RQ1).} HARNESS-DB releases every coded value with its evidence,
  at mean per-dimension Cohen's kappa 0.784 [0.765, 0.800] over 247 double-coded systems (37 of 38
  dimensions at or above 0.6; four intervals include the floor), which measures reproducibility, not
  correctness.

  \item \textbf{A weighting inversion (RQ2).} Repository-primary systems are 69.2\% of 1,252 coded
  systems carrying a value but an estimated 33.3\% of the frame, so a sample of well-known systems
  reverses the modal release artifact.

  \item \textbf{An under-reporting gradient (RQ4).} The weighted \texttt{not\_reported} rate runs from
  24.9\% $\pm$ 1.2 in the control-loop layer to 81.6\% $\pm$ 1.8 in the sandbox layer, a statement
  about documents, not about systems.

  \item \textbf{A measured trap (RQ2).} Mean corrected Cram\'er's V is 0.167, and treating silence as a
  category raises significant pairs from 385 to 648 of 666, 265 of them significant only when silence
  is a level.

  \item \textbf{A registered negative and a diversification (RQ2).} Clustering fails its silhouette
  floor, 0.120 against 0.25 (verdict negative); in the coded set, 10 of 34 dimensions show a diversity
  change whose interval
  excludes zero against 1.7 expected under the global null, of which 9 survive Benjamini--Hochberg
  (6 diversified, 3 converged).

  \item \textbf{A comparability ceiling and a selection-sensitive association (RQ3).} With 53 of 1,256 systems
  comparable to any peer, the registered regression was not fitted, and multi-agent topology is associated
  with $+0.88$ within-key standard deviations [$+0.38$, $+1.31$]. Drift attenuates it, but the choice of
  reference systems can inflate it: on the \rqXPOtherKeys{} keys whose references are not all recurring
  baselines the estimate is \rqXPOtherEffect{} \rqXPOtherCI{}, so the net direction of bias is unknown.

  \item \textbf{An upper bound that does not discriminate (RQ3).} On the completed corpus-wide harvest of
  authors' own ablations, which read \rqHarvestTierCoverage, \rqDiscountSurviveZ{} of
  \rqPoolableDims{} poolable dimensions pass the bias-discount rule with pooled relative gains of
  \rqEffectRangeSurviving{} (middle half \rqEffectIQRSurviving) before a median discount of \rqMedianDiscountPct\%, each an upper bound,
  and \rqSignSharePct\% of contrasts favour the ablated component: these ablations barely separate design
  dimensions. Self-verification, which the coded-set harvest singled out, is worth at most about
  \rqSVDiscounted{} relative; our compute-matched test of it failed its pilot band, and in 56
  diagnostic runs model-written self-checks accepted the first attempt 50 times, and 41 of 48 in a
  second pilot.

  \item \textbf{An evidence gap (RQ3).} On the coded-set harvest none of the eleven dimensions of
  layers F, G, and H carried a published ablation. The corpus-wide harvest reaches four of them,
  through 1 to 21 papers each; of the \rqPreRuleGContrasts{} contrasts it first attributed to layer G,
  \rqPreRuleGRemapped{} removed code-execution feedback or file-system tools and the other
  \rqPreRuleGDropped{} were unreadable or named no sandbox mechanism, and under the fixed mapping rules \rqZeroDimCount{} of
  \rqAblatableDimCount{} ablatable dimensions, all four of layer G among them, carry none. Ablation
  density falls as silence rises (Spearman's $\rho$ of \rqCoverageRho, exact one-sided permutation $p$
  of \rqCoverageP), an absence of evidence, not evidence that those components do not matter.
\end{enumerate}
